% Chapter V, Sections 53--54 and Bibliographical Notes
% Translated from: Chandrasekhar, "Hydrodynamic and Hydromagnetic Stability"
% Sections on thermal instability of a layer heated from below in the
% presence of rotation and magnetic field.

\section{The case when instability sets in as overstability}\label{sec:53}

We have seen that when liquid metals, such as mercury, are in rotation,
thermal instability sets in mostly as overstability. But in the presence of
a magnetic field, the principle of the exchange of stabilities is valid and
instability sets in as stationary convection. When rotation and magnetic
field are simultaneously present, the manner of the onset of instability
must, in general, depend in an extremely complicated way on the
relevant parameters $T$, $Q$, $p_1$ ($= \nu/\kappa$), and $p_2$ ($= \nu/\eta$). We shall seek
to unravel the nature of this dependence in the particular case of mercury.

Returning to equation~(43), and considering again the case of two free
boundaries adjoining a non-conducting medium, we can obtain the
required characteristic equation by inserting the solution $W = W_0 \sin \pi z$;
thus
\begin{align}
Ra^2(n^2{+}a^2{+}p_2\,\sigma)\bigl[(n^2{+}a^2{+}\sigma)(n^2{+}a^2{+}p_2\,\sigma){+}Qn^2\bigr]
\notag\\
= (n^2{+}a^2{+}p_1\,\sigma)\bigl\{(n^2{+}a^2)\bigl[(n^2{+}a^2{+}\sigma)(n^2{+}a^2{+}p_2\,\sigma){+}Qn^2\bigr]^2 +{} \notag\\
\qquad {}+ Tn^2(n^2{+}a^2{+}p_2\,\sigma)^2\bigr\},
\label{eq:5-62}
\end{align}
where it must be remembered that $\sigma$ can be complex. Letting
\begin{equation}
x = \frac{a^2}{\pi^2},\quad
i\sigma_i = \frac{\sigma}{\pi^2},\quad
R_* = \frac{R}{\pi^4},\quad
T_* = \frac{T}{\pi^4},\quad\text{and}\quad
Q_* = \frac{Q}{\pi^2},
\label{eq:5-63}
\end{equation}
we can rewrite equation~\eqref{eq:5-62} in the form
\begin{align}
R_* &= \frac{1+x}{x}\bigl(1{+}x{+}ip_1\,\sigma_i\bigr)\bigl(1{+}x{+}i\sigma_i\bigr)
+ Q_*\,\frac{1{+}x{+}ip_1\,\sigma_i}{1{+}x{+}ip_2\,\sigma_i}
\notag\\
&\quad{}+ \frac{T_*}{1{+}x}\;\frac{(1{+}x{+}ip_1\,\sigma_i)(1{+}x{+}ip_2\,\sigma_i)}
{\bigl(1{+}x{+}i\sigma_i\bigr)\bigl(1{+}x{+}ip_2\,\sigma_i\bigr)+Q_*}\,.
\label{eq:5-64}
\end{align}

Since our present interest is to specify the critical Rayleigh number
for the onset of instability via a state of purely oscillatory motions, it
will suffice to seek the conditions that equation~\eqref{eq:5-64} will allow solutions
for which $\sigma_i$ is real. Assuming then that this is the case, and equating
the real and the imaginary parts of equation~\eqref{eq:5-64}, we obtain the following
pair of equations for determining $R_*$ and $\sigma_i$ for given $x$, $Q_*$, and~$T_*$:
\begin{align}
R_* &= \frac{1+x}{x}\bigl[(1{+}x)^2 - p_1\,\sigma_i^2\bigr]
+ Q_*\,\frac{(1{+}x)^2{+}p_1\,p_2\,\sigma_i^2}{(1{+}x)^2{+}p_2^2\,\sigma_i^2}
\notag\\[6pt]
&\quad{}+ \frac{T_*}{1{+}x}\;\frac{\bigl[(1{+}x)^2{-}p_1\,\sigma_i^2{+}Q_*\bigr]\bigl[(1{+}x)^2{-}p_2\,p_1\,\sigma_i^2\bigr]
{+}(1{+}x)^2(p_1{+}p_2)(1{+}p_2)\sigma_i^2}
{\bigl[(1{+}x)^2{-}p_2\,\sigma_i^2{+}Q_*\bigr]^2{+}(1{+}x)^2(1{+}p_2)^2\sigma_i^4}
\label{eq:5-65}
\end{align}
and
\begin{align}
&\frac{T_*}{1{+}x}\;\frac{(p_1{-}1)(1{+}x)^2{+}(p_1{+}p_2)Q_*{+}p_2^2(p_1{-}1)\sigma_i^2}
{\bigl[(1{+}x)^2{-}p_2\,\sigma_i^2{+}Q_*\bigr]^2{+}(1{+}x)^2(1{+}p_2)^2\sigma_i^4}
\notag\\[6pt]
&\qquad{}+ Q_*\,\frac{p_1{-}p_2}{(1{+}x)^2{+}p_2^2\,\sigma_i^2}
+ 1 + p_1 = 0.
\label{eq:5-66}
\end{align}

For assigned values of $Q_*$ and $T_*$, equations~\eqref{eq:5-65} and~\eqref{eq:5-66} define $R_*$
as a function of $x$; the minimum of this function determines the critical
Rayleigh number for the onset of overstability (for the assigned $Q_*$ and
$T_*$). This minimum value of the Rayleigh number for the onset of over\-stability should be compared with the corresponding value for the onset
of convection as given by the formulae of \S\,52; the manner in which
instability will first set in will depend on which of the two Rayleigh
numbers is the smaller.

Of course, for certain ranges of the parameters, equation~\eqref{eq:5-66} will
not allow positive solutions for $\sigma_i^2$; for those ranges of the parameters,
overstability is not possible.

\subsection*{(a) An approximate solution applicable to liquid metals}

The equations giving $R_*$ and $\sigma_i$ simplify considerably when applied
to liquid metals such as mercury. These substances are characterized by
values of $p_2$ which are exceedingly small; thus for mercury at room
temperatures,
\begin{equation}
p_1 = \frac{1{\cdot}11}{4{\cdot}45}\times 10^{-1} = 0{\cdot}025;\qquad
p_2 = \frac{1{\cdot}11 \times 10^{-3}}{7{\cdot}6 \times 10^{4}} = 1{\cdot}5\times 10^{-7}.
\label{eq:5-67}
\end{equation}
Accordingly, in these cases, $p_2$ can be neglected in comparison with $p_1$
or unity. Moreover, it appears on examination that the terms in $p_2$ can
be neglected even when they occur multiplied by a large factor such as
$\sigma_i^2$. Therefore, suppressing all terms in equations~\eqref{eq:5-65}
and~\eqref{eq:5-66} which have $p_2$ as a factor, we obtain
\begin{equation}
R_* = \frac{1+x}{x}\bigl[(1{+}x)^2 - p_1\,\sigma_i^2\bigr]
+ Q_* + \frac{T_*}{1{+}x}\;\frac{(1{+}x)^2\bigl[(1{+}x)^2{+}Q_*\bigr]{+}\sigma_i^2}
{\bigl[(1{+}x)^2{+}Q_*\bigr]^2{+}(1{+}x)^2\sigma_i^2}\,,
\label{eq:5-68}
\end{equation}
and
\begin{equation}
(1{+}x)(1{+}p_1){+}Q_*\,\frac{p_1}{1{+}x} = \frac{T_*}{1{+}x}\;
\frac{(1{-}p_1)(1{+}x)^2 - p_1\,Q_*}
{\bigl[(1{+}x)^2{+}Q_*\bigr]^2{+}(1{+}x)^2\sigma_i^2}\,.
\label{eq:5-69}
\end{equation}
Equation~\eqref{eq:5-69} provides the following explicit formula for $\sigma_i^2$:
\begin{equation}
\sigma_i^2 = \frac{T_*}{1{+}x}\;
\frac{(1{+}x)^2(1{-}p_1) - p_1\,Q_*}{(1{+}x)(1{+}p_1)+p_1\,Q_*}
- \biggl[\frac{(1{+}x)^2{+}Q_*}{1{+}x}\biggr]^2.
\label{eq:5-70}
\end{equation}
Also, equations~\eqref{eq:5-68} and~\eqref{eq:5-69} can be combined to give
\begin{equation}
R_* = 2\,\frac{1+x}{x}\;\frac{(1{+}x)^2{+}Q_*}
{(1{+}x)^2(1{-}p_1)}\;\bigl[(1{+}x)^2{+}p_1\sigma_i^2\bigr].
\label{eq:5-71}
\end{equation}

\subsection*{(b) Numerical results for mercury}

Equations~\eqref{eq:5-70} and~\eqref{eq:5-71} have been used to determine the critical
Rayleigh numbers for the onset of overstability for various values of
$Q_*$ and $T_*$ for $p_1 = 0{\cdot}025$. The results of the calculations are summarized
in Table~\ref{tab:XIX}. The last column in this table gives the gyration frequency
of the oscillation in the marginal state in units of $\Omega$; according to equations~\eqref{eq:5-63} and the definition of~$T_*$:
\begin{equation}
p/\Omega = 2\sigma_i/\!\sqrt{T_*}\,.
\label{eq:5-72}
\end{equation}

\begin{table}[htbp]
\centering
\caption{Critical Rayleigh numbers and related constants for the onset of overstability
for the case $p_1 = 0{\cdot}025$}
\label{tab:XIX}
\medskip
{\small
\begin{tabular}{@{}r r r r r@{}}
\hline
$Q_*$ & $a_c$ & $R_{*c}$ & $R_c$ & $p/\Omega$ \\
\hline
\multicolumn{5}{c}{$T_* = 10^4$} \\[3pt]
10 & $4{\cdot}36$ & $33{\cdot}58$ & $7{\cdot}059\times10^3$ & $1{\cdot}0704$ \\
30 & $5{\cdot}35$ & $45{\cdot}35$ & $1{\cdot}388\times10^4$ & $0{\cdot}9110$ \\
100 & $6{\cdot}63$ & $90{\cdot}12$ & $3{\cdot}546\times10^4$ & $0{\cdot}6014$ \\
150 & $7{\cdot}30$ & $18{\cdot}44$ & $5{\cdot}019\times10^4$ & $0{\cdot}3688$ \\[6pt]
\multicolumn{5}{c}{$T_* = 10^5$} \\[3pt]
10 & $6{\cdot}03$ & $140{\cdot}8$ & $1{\cdot}284\times10^5$ & $0{\cdot}8905$ \\
20 & $6{\cdot}41$ & $139{\cdot}6$ & $1{\cdot}067\times10^5$ & $0{\cdot}8965$ \\
30 & $6{\cdot}66$ & $137{\cdot}8$ & $4{\cdot}039\times10^5$ & $0{\cdot}8803$ \\
40 & $6{\cdot}88$ & $133{\cdot}7$ & $2{\cdot}355\times10^4$ & $0{\cdot}7813$ \\
50 & $7{\cdot}03$ & $130{\cdot}4$ & $1{\cdot}499\times10^4$ & $0{\cdot}7615$ \\
70 & $7{\cdot}30$ & $117{\cdot}6$ & $1{\cdot}990\times10^4$ & $0{\cdot}7438$ \\
100 & $7{\cdot}66$ & $100{\cdot}9$ & $4{\cdot}317\times10^4$ & $0{\cdot}6894$ \\
200 & $8{\cdot}45$ & $94{\cdot}38$ & $7{\cdot}139\times10^4$ & $0{\cdot}5099$ \\
400 & $9{\cdot}48$ & $74{\cdot}53$ & $1{\cdot}278\times10^5$ & $0{\cdot}4587$ \\
300 & $9{\cdot}88$ & $61{\cdot}95$ & $1{\cdot}557\times10^5$ & $0{\cdot}3918$ \\[6pt]
\multicolumn{5}{c}{$T_* = 10^6$} \\[3pt]
10 & $8{\cdot}35$ & $341{\cdot}1$ & $4{\cdot}690\times10^5$ & $0{\cdot}6842$ \\
30 & $8{\cdot}89$ & $311{\cdot}7$ & $4{\cdot}533\times10^5$ & $0{\cdot}6434$ \\
100 & $9{\cdot}78$ & $293{\cdot}6$ & $7{\cdot}073\times10^5$ & $0{\cdot}6520$ \\
300 & $11{\cdot}34$ & $196{\cdot}5$ & $1{\cdot}893\times10^5$ & $0{\cdot}4760$ \\
1{,}000 & $12{\cdot}64$ & $202{\cdot}3$ & $3{\cdot}371\times10^5$ & $0{\cdot}4050$ \\
1{,}600 & $13{\cdot}62$ & $179{\cdot}3$ & $4{\cdot}898\times10^5$ & $0{\cdot}3406$ \\[6pt]
\multicolumn{5}{c}{$T_* = 10^7$} \\[3pt]
10 & $17{\cdot}4$ & $1739$ & $6{\cdot}000\times10^6$ & $0{\cdot}3458$ \\
100 & $18{\cdot}6$ & $1674$ & $6{\cdot}420\times10^6$ & $0{\cdot}3348$ \\
1{,}000 & $20{\cdot}9$ & $1407$ & $1{\cdot}002\times10^6$ & $0{\cdot}3136$ \\
10{,}000 & $25{\cdot}9$ & $1133$ & $3{\cdot}642\times10^6$ & $0{\cdot}2246$ \\
15{,}000 & $28{\cdot}4$ & $1024$ & $5{\cdot}348\times10^6$ & $0{\cdot}2048$ \\[6pt]
\multicolumn{5}{c}{$T_* = 10^8$} \\[3pt]
100 & $37{\cdot}7$ & $8097$ & $1{\cdot}354\times10^7$ & $0{\cdot}1619$ \\
1{,}000 & $38{\cdot}4$ & $7933$ & $1{\cdot}305\times10^7$ & $0{\cdot}1591$ \\
10{,}000 & $41{\cdot}9$ & $7289$ & $1{\cdot}720\times10^7$ & $0{\cdot}1446$ \\
40{,}000 & $46{\cdot}5$ & $6453$ & $2{\cdot}786\times10^7$ & $0{\cdot}1491$ \\
100{,}000 & $50{\cdot}1$ & $5841$ & $4{\cdot}376\times10^7$ & $0{\cdot}1168$ \\
130{,}000 & $51{\cdot}3$ & $5643$ & $5{\cdot}421\times10^7$ & $0{\cdot}1119$ \\
\hline
\end{tabular}
}
\end{table}

The critical Rayleigh numbers for the onset of overstability and
stationary convection are exhibited in Fig.~47 as functions of $Q_*$ for
various assigned values of $T_*$; the wave numbers of the associated disturbances are similarly exhibited in Fig.~48. We can directly read from
Fig.~47 the manner of the onset of instability for the various values of
$T_*$ for which the calculations have been made.

The phenomenon described in \S\,52 which gives rise to the peculiar,
non-monotone behaviour of the $(R, Q)$-curves when instability sets in
as stationary convection does not occur when allowance is made for the
onset of overstability for a value of $p_1$ as small as $0{\cdot}025$. However, the
curves do show a discontinuous behaviour though of a much less pro\-nounced character. From Fig.~47 it appears that for the cases illustrated,
the transition from overstability to convection occurs at about the place
where the convection curve passes through its minimum; and from Fig.~48 it is seen that this transition is accompanied by a substantial dis\-continuity in the wave number of the disturbance which manifests itself
at marginal stability. This discontinuity is in the sense that the convection cells (for increasing $Q_*$) suddenly get very much widened.

%% ==================================================================
\section{Experiments on the onset of thermal instability in the
presence of rotation and magnetic field}\label{sec:54}

Experiments on the onset of thermal instability in the presence of
rotation and magnetic field have been successfully carried out by
Nakagawa.

In Nakagawa's experiments, a reconditioned electromagnet of a $32\tfrac{1}{2}$-in.\
cyclotron of the University of Chicago was used to provide a uniform
vertical magnetic field of variable strength. A Pyrex glass cylinder
containing mercury mounted on levelled non-magnetic ball-bearings
between the pole pieces of the electromagnet was rotated in the magnetic
field by a constant speed electric motor (see Fig.~50 which shows the
arrangement used in the experiments described in \S\,(b) below). The
angular velocity of the rotation was measured by an optical method using
electrical timing devices; the accuracy attained in the time measurements was $10^{-3}$~sec.

A direct-current electric heater placed at the bottom of the container,
below a stainless-steel plate, provided the means of heating. The
temperature gradient which was established was measured by a nine-couple copper--constantan thermopile immersed at two fixed levels in
the mercury. The average temperature of the mercury in any particular
experiment was also measured by a thermocouple, one of whose junctions
was placed at the middle of the mercury layer; the values of the various
coefficients (like $\nu$, $\kappa$, etc.)\ at this average temperature were used in
evaluating the parameters such as $R$, $T$, and~$Q$.

The e.m.f.\ of the thermopile and the thermocouple were amplified
and recorded automatically and continuously. The accuracy in the
thermopile measurements was $\pm 0{\cdot}001^\circ$~C, while the average temperatures measured by the thermocouple were accurate to $\pm 0{\cdot}01^\circ$~C.

The electrical connexions to the various components in the rotating
part of the assembly were made through a system of mercury filled
troughs with copper and constantan contacts.

A constant circulation of cooled nitrogen helped to maintain the
conditions at the top surface of the mercury constant and uniform.

\figurePlaceholder{50}{A schematic diagram of the experimental arrangement:
$A$, bakelite cylinder; $B$, stainless-steel plate; $C$, electric heater;
$D$, non-magnetic ball-bearing; $E$, stainless-steel rod; $F$, mercury trough;
$M$, front-surface mirror; $S$, rotary shutter; $T$, camera
(\textit{Proc.\ Roy.\ Soc.\ (London)} A, \textbf{249}, 140 (1958)).}

\subsection*{(a) The results on the critical Rayleigh number and on the manner of
the onset of instability}

All the experiments were carried out with a layer of mercury 3~cm in
depth rotated with a constant angular velocity of 5~rev/min. The corresponding value of $T_*$ varied between $7{\cdot}5\times 10^4$ and $8{\cdot}5\times 10^6$ depending on
the average temperature of the mercury. Experiments were performed
for values of $Q_*$ in the range 10 to $2\times 10^4$ by using magnetic fields of
strength varying from 125 to 5{,}000~gauss.

The critical Rayleigh number for the onset of instability, for a constant
speed of rotation and strength of magnetic field, was determined by
measuring the steady temperature gradients which get established for
various rates of heating, and using the Schmidt--Milverton principle;
and the manner of the onset of instability was ascertained in each case
by inspecting the temperature records. The onset of overstability could
always be distinguished by the nearly pure sinusoidal oscillations
exhibited by the temperature records. A typical example of such a
record is shown in Fig.~51.

\figurePlaceholder{51}{A time record of the adverse temperature gradient for
mercury; $d = 3$~cm; $\Omega = 5$~rev/min; $H = 125$~G; $Q_* = 1{\cdot}01\times 10^1$;
$T = 7{\cdot}90\times 10^4$.}

In Fig.~52 the results of the experiments together with the theoretically
expected relation for $T_* = 10^6$, according to the calculations of \S\S\,52
and~53, are shown. Since the theoretical results have been derived for two
free boundaries and refer, moreover, to a value of $T_*$ somewhat different
from the average value to which the experiments correspond, a quantitative agreement between the experiments and the theory is not to be
expected. All that can be expected is an agreement in the general
features; and this is certainly present. In particular, the characteristic
discontinuity in the dependence of the Rayleigh number on $Q_*$ with the
accompanying change in the manner of the onset of instability is clearly
demonstrated.

The periods $2\pi/p$ of the overstable oscillations were estimated by
counting the number of oscillations in the temperature records which
occur in known intervals of time. In Fig.~53 these results are shown
together with the theoretically expected relations for $T_* = 10^5$ and~$10^6$.
Again the general agreement is satisfactory, though as we have explained,
quantitative agreement is not to be expected.

\figurePlaceholder{52}{A summary of the experimental and the theoretical results on the
critical Rayleigh numbers for instability. The curves labelled $T_* = 10^6$ (convection), $T_* = 10^6$ (overstability), and $T_* = 0$ (convection) represent the
theoretically derived relations. The value of $T_*$ appropriate for the experimental points represented by the solid and shaded triangles is $7{\cdot}78\times 10^5$.}

\figurePlaceholder{53}{A comparison of the observed periods of overstable oscillations with
the theoretical values (the full line curves). The value of $T_*$ appropriate for
the experimental results is $(6{\cdot}05\pm 0{\cdot}07)\times 10^5$.}

\subsection*{(b) Optical observations and the discontinuous variation of the cell
dimensions with the strength of the magnetic field}

For optical observations, some slight modifications in the arrangement
described earlier were necessary. It is the arrangement used in these
experiments that is illustrated in Fig.~50.

Mercury was contained in a Bakelite cylinder 24~cm in diameter and
4~cm in height. The bottom of the cylinder $A$ was sealed by a stainless-steel plate $B$ with an O-ring. The steel plate was carefully machined to
eliminate surface irregularities which might disturb the ensuing convection patterns; freedom from such disturbances was found to be essential
for successful optical observations.

Uniform heating from below was provided by a non-inductively
wound electric heater $C$ placed below the stainless-steel plate, and
the electric heater was mounted, between the pole pieces of the electromagnet, on non-magnetic ball-bearing~$D$. And this assembly was rotated
in the magnetic field at 5~rev/min as in the previous experiments.

The optical part of the apparatus consisted of a mirror $M$ mounted $45^\circ$
to the vertical, the rotoscope, a rotary shutter~$S$, and a single lens reflex
camera~$T$. The mirror was silvered on the front surface; and dimming
of the mirror (by condensation of water-vapour during the experiments)
was prevented by heat from an infra-red lamp.

A stationary image was secured at the camera by rotating the rotoscope with an angular velocity exactly one-half that of the apparatus.
By a suitable sequence of openings and closings of the shutter, it was
possible to obtain photographs every 15~sec.

As tracers of the motion, sand particles of diameter approximately
$\frac{1}{4}$~mm were used. These particles floating freely on the surface, and
following the motion of the mercury, enabled the emerging pattern of
convection to be photographed. A thin layer of distilled water on the
mercury surface eliminated much of the difficulty arising from surface
oxidation.

The experiments were again carried out on a layer of mercury 3~cm in
depth and an angular velocity of rotation of 5~rev/min. And as ascertained by the experiments in~\S\,(a), the rate of heating for each value of $Q_*$
was adjusted so that the conditions were just marginal. By repeating
the experiments at least twice for each value of $Q_*$, Nakagawa took a
total of 300 streak photographs. Measurements were made on prints of
these photographs on $8\times 10$-in.\ paper. Distances between the centres
of adjacent cells were measured and in the manner described in Chapter~IV, \S\,48, the theoretical values to which these distances should correspond were inferred. In Fig.~54 the results of the measurements are
compared with the expected theoretical relation for $T_* = 10^6$. While,
for reasons we have already explained, a quantitative agreement is not
to be expected, the agreement in the general features of the relationship

\figurePlaceholder{54}{A summary of the experimental and theoretical results on the sizes
of the cells manifested at marginal stability. The theoretical relations for
$T_* = 10^6$ and $T_* = 0$ are represented by the broken lines. Solid circles
represent the experimental results for $T_* = 7{\cdot}30\times 10^5$; and the heavy solid lines
are fitted to the experimental data. The lower branches of the broken line
and the heavy solid line represent the case when the instability sets in as
overstability. Both lines on the right-hand side of the figure represent the
case when the instability sets in as steady convection.}

\noindent
is satisfactory. In particular, the predicted enlargement of the cell
dimension at a critical field strength is strikingly confirmed.

In Fig.~55 some typical examples from Nakagawa's collection of
streak photographs are illustrated. They illustrate the characteristic
dependence of the size of the cells on the value of $Q_*$. The discontinuous
change in the size of the cells which occurs for increasing field strength is
clearly demonstrated by this sequence of photographs.

\figurePlaceholder{55}{Examples of streak photographs of the convective motion on the surface of a layer
of mercury of depth 3~cm, rotating at 5~rev/min, and subject to impressed external magnetic
fields of various strengths: (a) $H = 125$~G, $Q_* = 9{\cdot}60$, and $T_* = 6{\cdot}95\cdot 10^5$;
(b) $H = 750$~G, $Q_* = 3{\cdot}52\cdot 10^3$, and $T_* = 7{\cdot}30\cdot 10^5$;
(c) $H = 1{,}000$~G, $Q_* = 6{\cdot}25\cdot 10^3$, and $T_* = 7{\cdot}31\cdot 10^5$;
(d) $H = 3{,}000$~G, $Q_* = 5{\cdot}61\cdot 10^4$, and $T_* = 7{\cdot}49\cdot 10^5$.
The sudden enlargement of the cell size as we pass from $H = 750$~G to $H = 1{,}000$~G is apparent.}

%% ==================================================================
\section*{BIBLIOGRAPHICAL NOTES}
\addcontentsline{toc}{section}{Bibliographical Notes}

\paragraph{\S\,50.} Hydromagnetic waves in a rotating fluid have been considered by:
\begin{enumerate}
\item B.~\textsc{Lehnert}, `Magnetohydrodynamic waves under the action of the
Coriolis force, Part~I', \textit{Astrophys.\ J.}\ \textbf{119}, 647--54 (1954);
\textbf{Part~II}, \textit{ibid.}\ \textbf{121}, 481--90 (1955).
\end{enumerate}
See also:
\begin{enumerate}\setcounter{enumi}{1}
\item S.~\textsc{Chandrasekhar}, `The gravitational instability of an infinite homogeneous medium when Coriolis force is acting and a magnetic field is present',
\textit{ibid.}\ \textbf{119}, 7--9 (1954).
\item B.~\textsc{Lehnert}, `The decay of magneto-turbulence in the presence of a
magnetic field and Coriolis force', \textit{Quart.\ Appl.\ Math.}\ \textbf{12}, 321--41 (1955).
\end{enumerate}

\paragraph{\S\S\,51--53.} The analyses in these sections are based on:
\begin{enumerate}\setcounter{enumi}{3}
\item S.~\textsc{Chandrasekhar}, `The instability of a layer of fluid heated below and
subject to the simultaneous action of a magnetic field and rotation.~I',
\textit{Proc.\ Roy.\ Soc.\ (London)}~A, \textbf{225}, 173--84 (1954); II, \textit{ibid.}\ \textbf{237}, 476--84 (1956).
\end{enumerate}

The fact that the curve $R(x)$ defined by equations~(59) and~(60) has two minima
for certain ranges of the parameters $Q$ and $T$, was first observed by Donna Elbert.

\paragraph{\S\,54.} The references for the experimental work described in this section are:
\begin{enumerate}\setcounter{enumi}{4}
\item Y.~\textsc{Nakagawa}, `Experiments on the instability of a layer of mercury
heated from below and subject to the simultaneous action of a magnetic
field and rotation.~I', \textit{Proc.\ Roy.\ Soc.\ (London)}~A, \textbf{242}, 81--88 (1957);
II, \textit{ibid.}\ \textbf{249}, 138--45 (1959).
\end{enumerate}

Extending the experiments described in reference~5, Nakagawa has since demonstrated the occurrence of instability as stationary convection after the instability
has already manifested itself as overstability; and he has, moreover, observed
the predicted decrease of the Rayleigh number, with increasing strength of the
magnetic field, over a part of the convective branch (see Fig.~52). [For the corresponding results for the case when rotation alone is present, see pp.~142 and~143.]
